\par\Needspace{4\baselineskip}
\begin{center}
\hypertarget{native-042-article025}{}\pdfbookmark[2]{25.}{article-025}25.
\end{center}
Beim inversen Problem sollen aus der Zeit die wahre Anomalie und der Radiusvektor bestimmt werden. Dazu ermittelt man zunächst aus $N=\lambda kb^{-\frac{3}{2}}t$ mithilfe von Gleichung XI die Hilfsgröße $u$ oder $F$. Diese transzendente Gleichung ist durch Probieren zu lösen. Mit Verfahren, die den in Art. 11 erläuterten ähneln, lässt sich die Rechnung abkürzen. Auf eine ausführlichere Erläuterung verzichten wir. Hyperbolische Bewegung dürfte am Himmel kaum je vorkommen, sodass es sich wohl nicht lohnt, die Rechenvorschriften dafür ebenso sorgfältig auszuarbeiten wie für die elliptische Bewegung. \SourceLeaf{043}Außerdem lassen sich alle etwa auftretenden Fälle mit einem anderen Verfahren behandeln, das wir weiter unten angeben. Hat man $F$ oder $u$ gefunden, so bestimmt man daraus zunächst $v$ nach Formel III und anschließend $r$ nach Formel II oder VIII. Noch bequemer lassen sich $v$ und $r$ zugleich nach den Formeln VI und VII ermitteln. Zur Kontrolle der Rechnung kann man nach Belieben die eine oder andere der übrigen Formeln heranziehen.

\par\Needspace{4\baselineskip}
\begin{center}
\hypertarget{native-043-article026}{}\pdfbookmark[2]{26.}{article-026}26.
\end{center}
Beispiel. $e$ und $b$ bleiben wie im vorigen Beispiel. Es sei $t=65{,}41236$. Gesucht sind $v$ und $r$. Mit den \textsc{Briggs}schen Logarithmen erhält man
\begin{center}
\begin{tabular}{@{}ll@{}}
\makebox[18em][l]{$\log t$\hfill .\,.\,.\,.\,.\,.\,.\,.\,.\,.\,.\,.} & $1{,}815\,6598$\\[.3em]
\makebox[18em][l]{$\log\lambda kb^{-\frac{3}{2}}$\hfill .\,.\,.\,.\,.\,.\,.\,.} & $6{,}970\,2758$\\[.3em]
\hline
\makebox[18em][l]{$\log N$\hfill .\,.\,.\,.\,.\,.\,.\,.\,.\,.\,.} & $8{,}785\,9356$,
\end{tabular}\\[.25em]
woraus $N=0{,}061\,0851\,4$ folgt.
\end{center}
Daraus ergibt sich für die Gleichung $N=\lambda e\operatorname{tang}F-\log\operatorname{tang}(45^0+\frac{1}{2}F)$ der folgende Wert als Lösung
\[
F=25^0\,24^{\prime}\,27^{\prime\prime}{,}66,
\]
woraus sich nach Formel III ergibt
\begin{center}
\begin{tabular}{@{}ll@{}}
\makebox[18em][l]{$\log\operatorname{tang}\frac{1}{2}F$\hfill .\,.\,.\,.\,.\,.\,.} & $9{,}353\,0120$\\[.3em]
\makebox[18em][l]{$\log\operatorname{tang}\frac{1}{2}\psi$\hfill .\,.\,.\,.\,.\,.\,.} & $9{,}531\,8179$\\[.3em]
\hline
\makebox[18em][l]{$\log\operatorname{tang}\frac{1}{2}v$\hfill .\,.\,.\,.\,.\,.\,.} & $9{,}821\,1941$
\end{tabular}
\end{center}
also $\frac{1}{2}v=33^0\,31^{\prime}\,29^{\prime\prime}{,}89$ und $v=67^0\,2^{\prime}\,59^{\prime\prime}{,}78$. Ferner erhält man
\begin{center}
\footnotesize
\begin{tabular}{@{}ll@{}}
$\left.\begin{array}{ll}
\text{C.}\log\cos\frac{1}{2}(v+\psi)\quad .\,.\,. &0{,}213\,7476\\[.3em]
\text{C.}\log\cos\frac{1}{2}(v-\psi)\quad .\,.\,. &0{,}014\,5197
\end{array}\right\}$ & \begin{tabular}{@{}ll@{}}
Differenz .\,.\,.\,.\,.\,.\,.\,.\,.\,.\,.\,. & $0{,}199\,2279$\\[.3em]
$\log\operatorname{tang}(45^0+\frac{1}{2}F)$ .\,.\,. & $0{,}199\,2280$
\end{tabular}
\end{tabular}
\end{center}
\begin{center}
\begin{tabular}{@{}ll@{}}
\makebox[18em][l]{$\log\frac{p}{2e}$\hfill .\,.\,.\,.\,.\,.\,.\,.\,.\,.\,.} & $9{,}972\,5868$\\[.3em]
\hline
\makebox[18em][l]{$\log r$\hfill .\,.\,.\,.\,.\,.\,.\,.\,.\,.\,.\,.} & $0{,}200\,8541$
\end{tabular}
\end{center}
